\section*{Chapitre \ref{ch:s1:merger-arbitrage} --- Arbitrage de fusions}

\begin{solution}{exo:s1:merger-arbitrage:1}
\omterm{def:s1:merger-arbitrage:spread}{Spread de l'opération} $40/38.60 - 1 = 3.63\%$~; probabilité implicite $(38.60 - 30)/(40 - 30) = 0.86$.
\end{solution}

\begin{solution}{exo:s1:merger-arbitrage:2}
$(1.0363)^{252/84} - 1 = 11.3~\%$ par an.
\end{solution}

\begin{solution}{exo:s1:merger-arbitrage:3}
À 42~\$, dans la bande, la contrepartie vaut $0.5 \times 42 = \$21$ et la couverture consiste à vendre à découvert 0.5 action de l'acquéreur. À 36~\$, sous la bande, la contrepartie est
fixée à $0.5 \times 40 = \$20$ et la couverture est nulle~: la valeur ne dépend plus de l'acquéreur.
\end{solution}

\begin{solution}{exo:s1:merger-arbitrage:4}
Profit espéré $0.86 \times (40 - 38.60) - 0.14 \times (38.60 - 30) = 1.204 - 1.204 = 0$. Pour prendre position, l'arbitragiste doit croire que la probabilité d'échec est inférieure à
14~\% (ou le \omterm{def:s1:merger-arbitrage:spread}{prix en cas d'échec} supérieur à 30~\$), après actualisation et frais.
\end{solution}

\begin{solution}{exo:s1:merger-arbitrage:5}
Le \omterm{def:s1:merger-arbitrage:spread}{prix en cas d'échec} de chaque opération évolue avec le marché~: une baisse du marché fait baisser la valeur de toutes les positions ouvertes d'un coup, qu'une opération échoue ou non~; et davantage d'opérations échouent
lorsque le marché baisse. Ces deux effets sont communs à toutes les opérations et ne se diversifient pas.
\end{solution}

\begin{solution}{exo:s1:merger-arbitrage:6}
Elle allonge l'examen de plusieurs mois et signale une préoccupation antitrust~: la probabilité d'échec augmente et le délai de réalisation s'allonge, si bien que le \omterm{def:s1:merger-arbitrage:spread}{spread de l'opération} s'élargit et que son
rendement annualisé baisse pour un même spread en valeur absolue.
\end{solution}

\begin{solution}{exo:s1:merger-arbitrage:7}
Si les échecs sont indépendants du marché, le portefeuille rapporte 4.30~\% au-dessus du taux (ratio de Sharpe 1.50)~: la concentration des échecs lors des replis coûte 1.5 point par
an. Avec des spreads évalués au taux d'échec d'un marché plat, le portefeuille perd 0.47~\% (ratio de Sharpe $-0.17$)~: la prime du spread de base rémunère les
échecs liés aux krachs.
\end{solution}

\begin{solution}{exo:s1:merger-arbitrage:8}
Son bêta proche de zéro est une moyenne~: il est voisin de zéro sur les marchés haussiers et plats, et plus élevé sur les marchés baissiers (0.20 contre 0.09 ici, davantage dans les
données historiques), à la manière d'un put vendu. Il est vendeur de risque de krach, et un ratio de Sharpe construit sur des mois sans krach surestime sa qualité.
\end{solution}

\begin{solution}{pb:s1:merger-arbitrage:1}
\begin{enumerate}
\item Acheter les cibles en dessous de la contrepartie (couverture des offres en titres) pour gagner le spread à la réalisation~; une offre faite directement aux actionnaires~; offre rapportée au cours,
moins un~; le prix si l'opération échoue.
\item 0.86.
\item Les opérations durent plusieurs mois~; un même spread sur des horizons différents correspond à un rendement différent.
\item Environ cinquante offres en numéraire par an, des primes de 20--40~\%, une durée médiane de 110 jours, des \omterm{def:s1:merger-arbitrage:spread}{prix en cas d'échec} qui évoluent avec le marché, des spreads traduisant une probabilité
d'échec de 12~\%, probabilité réelle de 6~\% plus 1.5 fois la baisse du marché.
\item 6.76~\% par an, 2.76~\% au-dessus du taux, volatilité de 3.26~\%, ratio de Sharpe 0.85.
\item 88~\% implicite, 92.5~\% réalisée.
\item 13.7~\% lorsque le marché a baissé pendant la vie de l'opération, 4.7~\% lorsqu'il a monté.
\item 0.20 les mois où le marché a baissé de plus de 4~\%, 0.09 sinon.
\item Des échecs indépendants rapportent 4.30~\%~: la concentration coûte 1.5 point~; des spreads évalués à 6~\% perdent 0.47~\%~: la prime rémunère le risque de krach.
\item $-3.06~\%$ dans un mois où le marché a monté de 2.8~\%~: quelques échecs propres à certaines opérations survenant ensemble.
\item Vendre à découvert le ratio d'échange en actions de l'acquéreur~; avec un collar, vendre à découvert son delta, qui saute aux bords de la bande.
\item Examen antitrust, financement, votes des actionnaires, offres concurrentes, changements défavorables significatifs.
\item \textbf{Résultat nommé.} Le bêta du portefeuille vaut 0.20 les mois où le marché baisse de plus de 4~\% et 0.09 les autres mois~; les spreads impliquent une
probabilité de réalisation de 88~\%, et 92.5~\% des opérations aboutissent.
\item Rendements corrélés au marché lors des fortes baisses et non autrement, à la manière d'un put vendu sur indice~; environ 4~\% par an de rendement excédentaire après frais, 1963--1998.
\item Un put vendu sur indice hors de la monnaie.
\item Par la perte en cas d'échec, rapportée au portefeuille.
\item Au moins 30 jours (15 pour les \omterm{def:s1:merger-arbitrage:arb}{offres publiques} en numéraire) après le dépôt complet des dossiers, davantage en cas de Second Request.
\item Les échecs propres à l'opération, et non les échecs et \omterm{def:s1:merger-arbitrage:spread}{prix en cas d'échec} dépendant du marché.
\item Après une baisse, lorsque les spreads sont larges, si le capital peut supporter de nouvelles baisses.
\item Une assurance contre l'échec des opérations, qui rapporte le plus sur les marchés difficiles.
\end{enumerate}
\end{solution}

\begin{solution}{iq:s1:merger-arbitrage:1}
Acheter la cible d'une acquisition annoncée en dessous de l'offre (et vendre à découvert l'acquéreur dans les offres en titres) pour gagner le spread à la clôture de l'opération. Le rendement
est la prime versée pour porter le risque d'échec, dont une bonne part survient lors des replis de marché.
\end{solution}

\begin{solution}{iq:s1:merger-arbitrage:2}
Vendre à découvert le delta du collar en actions de l'acquéreur~: le ratio d'échange à l'intérieur de la bande, nul à l'extérieur où la valeur est fixée~; rééquilibrer à mesure que le cours de l'acquéreur
s'approche des bords de la bande, où le delta saute, ou acheter des options compensant le gamma du collar.
\end{solution}

\begin{solution}{iq:s1:merger-arbitrage:3}
Son profil de gains est une exposition proche de zéro sur les marchés normaux et des pertes lors des krachs, lorsque les opérations échouent et que les \omterm{def:s1:merger-arbitrage:spread}{prix en cas d'échec} baissent ensemble~: de petites primes régulières assorties de rares
pertes importantes et corrélées, soit le profil d'un put vendu.
\end{solution}

\begin{solution}{iq:s1:merger-arbitrage:4}
\omterm{def:s1:merger-arbitrage:spread}{Prix en cas d'échec} obtenu à partir du cours non perturbé, corrigé de l'évolution du marché et du secteur depuis l'annonce, ou à partir de pairs~; probabilité d'échec déduite des caractéristiques de l'opération
(recouvrement réglementaire, financement, hostilité, vote), taux de base d'opérations comparables, et probabilité implicite de marché comme a priori.
\end{solution}

\begin{solution}{iq:s1:merger-arbitrage:5}
Rumeurs de difficultés réglementaires ou de financement, une baisse du marché modifiant les \omterm{def:s1:merger-arbitrage:spread}{prix en cas d'échec}, ventes forcées d'autres arbitragistes, ou
problème d'emprunt de l'acquéreur dans une offre en titres.
\end{solution}

\begin{solution}{iq:s1:merger-arbitrage:6}
Avec le prix $A = e^{-r\tau}(pK + (1 - p)B)$, la probabilité implicite vaut $p = (Ae^{r\tau} - B)/(K - B)$. Si la probabilité réelle est $\pi$, le gain espéré est
$\pi K + (1 - \pi)B$ et le rendement excédentaire espéré vaut $(\pi - p)(K - B)/(Ae^{r\tau})$~: proportionnel à l'écart entre les probabilités réelle et implicite
et à l'écart entre l'offre et le \omterm{def:s1:merger-arbitrage:spread}{prix en cas d'échec}.
\end{solution}
